\section{The rational incidence cover}
\label{sec:cover}

\subsection{The Clebsch chart}

Put \(E=\Q[t]/(t^2-t-1)=\Q(\sqrt5)\), and let
\[
\begin{aligned}
 I_t=\{&
 [0:t:1],[0:t:-1],[1:0:t],[-1:0:t],\\
 & [t:-1:0],[-t:-1:0]\}\subset\mathbf P^2_E .
\end{aligned}
\]
These are the six axes of the standard icosahedron.  Define the linear
subspace
\[
 V_t=\{p\in H_E:p(a)=0\text{ for every }[a]\in I_t\}\subset H_E.
\]
This is the exact inclusion used below.  Hitchin proves that \(V_t\) has
dimension four and describes it as follows
\cite[\S\S2--4]{HitchinIcosahedron}.  The rotation group of \(I_t\) permutes
the five triples of mutually orthogonal symmetry planes.  Let
\(q_1,\ldots,q_5\) be the products of the three linear equations in those
triples, with \(q_1=xyz\) and relative scalars chosen so that
\(\sum_iq_i=0\).  Then
\[
 \iota_t:V_E=\{(y_i)\in E^5:\textstyle\sum_i y_i=0\}
   \longrightarrow V_t,\qquad
 (y_i)\longmapsto\sum_i y_iq_i
\]
is an \(A_5\)-equivariant isomorphism.  Thus the parameter of \(xyz=q_1\)
is
\[
 y^\circ=\frac15(4,-1,-1,-1,-1),\qquad
 \sigma_3(y^\circ)=\frac4{25}.
\]
We first make the sextic scale precise.  For a real harmonic cubic \(f\),
divide Hitchin's symmetric moment operator by \(4\pi\) and write
\[
 M_f=\frac1{4\pi}\int_{S^2}f(x)^2
       \left(xx^{\mathsf T}-\frac4{15}I\right)d\omega
\]
Hitchin applies \(3a_3-4a_1a_2\) to the coefficients of its characteristic
polynomial
\(\lambda^3+a_1\lambda^2+a_2\lambda+a_3\).  Normalized spherical moments
of monomials are rational, so this is a rational sextic polynomial on
\(H\).  Its zero set is Hitchin's sextic.  We record its irreducibility
paper-locally.  After complexification, let
\(S_{\mathbf C}\subset\mathbf P(V_{\mathbf C})\) be the Clebsch diagonal
cubic surface.  It is smooth, hence irreducible, and
\(\operatorname{SO}(3,\mathbf C)\) is connected, hence irreducible.  Therefore
the Zariski closure of the image of the action morphism
\[
 \operatorname{SO}(3,\mathbf C)\times S_{\mathbf C}
 \longrightarrow \mathbf P(H_{\mathbf C})
\]
is irreducible.  Hitchin's orbit description identifies this closure with
the sextic hypersurface \cite[\S9]{HitchinIcosahedron}.  Thus the hypersurface,
and hence its rational scalar normalization \(J_0=0\) below, is irreducible
over \(\mathbf C\).  In the appendix Hitchin rescales the moment operator
before the final restriction calculation; we therefore use the opening
formula only to fix the hypersurface, not to identify the two printed scales.

The pullback of any rational equation of this hypersurface is a nonzero
scalar multiple of \(\sigma_3^2\).  That scalar lies in \(\Q\): applying
\(t\mapsto1-t\) replaces the chart by its conjugate without changing either
the rational equation or \(\sigma_3^2\).  There is consequently a unique
rational multiple, denoted \(J_0\), for which Hitchin's appendix calculation
\cite[Appendix, final displayed formula]{HitchinIcosahedron} reads
\[
 \iota_t^*J_0=16\sigma_3^2,
 \qquad J_0(xyz)=\left(\frac{16}{25}\right)^2.
\]
This is an internal rational normalization, not an identification of two
analytic scales in the source.  The coefficient is fixed by the chosen
lift: at \(y^\circ\), the ratio of \(J_0(xyz)\) to
\(\sigma_3(y^\circ)^2\) is exactly \(16\).  Over \(\mathbf R\), \(J_0\) is a
positive multiple of Hitchin's opening invariant, so his sign
classification is unchanged.

This chart is defined over \(E\), not over \(\Q\).  The nontrivial
automorphism \(t\mapsto1-t\) carries \(I_t\) and \(V_t\) to the distinct
conjugate icosahedron and its vanishing space; the matrix \(R\) in
Section~\ref{sec:arithmetic} gives the corresponding transport.  Hence
\(\iota_t\) is not a literal rational inclusion \(V\hookrightarrow H\).
The two displayed icosahedra have no common axis.  Hitchin's
two-icosahedron intersection theorem therefore gives
\[
 V_t\cap V_{1-t}=E\cdot xyz.
\]
This is Hitchin's intersection theorem
\cite[Proposition~7]{HitchinIcosahedron}.
Consequently the reduced union
\(\mathbf P(V_t)\cup\mathbf P(V_{1-t})\) descends to \(\Q\): over \(E\) it
is a pair of conjugate projective three-spaces meeting only at the rational
point \([xyz]\).  Denote the descended scheme by \(Z\).  Its normalization
map is \(\nu:Z^\nu\to Z\), where \(Z^\nu=\mathbf P(V_t)\), regarded as a
\(\Q\)-scheme through
\(\Spec E\to\Spec\Q\); after base change to \(E\), this normalization is
the disjoint union of the two conjugate components.  The inverse image of
the unique intersection point is
\[
 \nu^{-1}([xyz])=\Spec E.
\]
The local algebra is an instance of a general pinching construction.

\smallskip
\noindent\emph{Quadratic pinching.}\par
Let \(E/k\) be a separable quadratic field extension, put
\(A=E[[x_1,\ldots,x_n]]\) with \(n\geq1\), and let \(\mathfrak m\) be its
maximal ideal.
The fibre product
\[
 R=A\mathbin{\times_E}k=k+\mathfrak m
\]
has normalization \(A\), conductor \(\mathfrak m\), and defect
\(A/R\simeq E/k\) as a \(k\)-vector space.  After completed base change to
\(E\), it is the union
of two formal \(n\)-planes meeting at their origins.

Choose \(\alpha\in E\setminus k\).  Then \(A=R+R\alpha\), so \(A\) is
finite, hence integral, over \(R\).  The two rings have the same fraction
field because \(\alpha=(\alpha x_1)/x_1\).  Since \(A\) is normal, it is
the normalization.  An element of the conductor first lies in \(R\), so its
residue is some \(c\in k\).  Multiplication by the constants in \(E\)
forces \(cE\subset k\), hence \(c=0\); conversely
\(\mathfrak m A\subset\mathfrak m\).  Thus the conductor is
\(\mathfrak m\), and reduction modulo \(\mathfrak m\) gives
\(A/R\simeq E/k\) as \(k\)-vector spaces.  Finally
\(E\otimes_kE\simeq E\times E\); the fibre-product description gives the
two conjugate branches after completed base change.

The two projective three-spaces above meet transversely, so the completed
local ring of \(Z\) is this pinching ring with \(k=\Q\),
\(E=\Q(\sqrt5)\), and \(n=3\).  Its normalization, conductor,
and defect \(E/\Q\) are therefore forced by residue-field pinching; they do
not depend on a six-variable presentation.  Thus \(5\) is the discriminant
square class of the two branches, not merely a value recovered from the
sextic.  This local model explains the descent geometry but does not replace
the comparison with Hitchin's incidence scheme below.
After base change to \(E\), the pullback of the incidence double cover to
\(D(\sigma_3)\) is split.  Galois conjugation exchanges both the two
components and the two conjugate charts.

\subsection{A rational incidence model}

The function-field calculation requires a rational model, which can be
given without descending either Clebsch chart.  Let \(W=\Q^3\) with the
quadratic form \(Q\) from the introduction.  For \(a\in W\), infinitesimal
rotation about \(a\) acts on \(H\); put
\[
 \omega_a(p,q)=\langle a\mathbin{\cdot}p,q\rangle .
\]
These three rational skew forms define a closed subscheme
\[
 X=\{U\in\operatorname{Gr}(3,H):
       \omega_a|_U=0\text{ for every }a\in W\}.
\]
Hitchin identifies \(X_{\mathbf C}\) with the Mukai--Umemura threefold
\cite[\S\S4--5]{HitchinBundles}; in particular \(X\) is geometrically
integral and smooth.  The incidence scheme used here is
\[
 \mathcal I=\{([f],U)\in\mathbf P(H)\times X:
                   \langle f,U\rangle=0\}.
\]
Let \(\pi:\mathcal I\to\mathbf P(H)\) be the first projection.
It is a projective bundle over \(X\), is therefore normal and integral,
and all its displayed equations are defined over \(\Q\).  Mukai and Umemura
introduced the underlying threefold; Hitchin supplies this Grassmannian
model and the incidence calculation \cite{MukaiUmemura,HitchinBundles}.

There is also a canonical link from the chart model to the finite
incidence cover.  Let
\[
 \mathcal N=\operatorname{Spec}_{\mathbf P(H)}\pi_*\mathcal O_{\mathcal I}.
\]
Stein factorization gives
\(\mathcal I\to\mathcal N\to\mathbf P(H)\).  Since \(\mathcal I\) is normal
and the degree calculation below makes its generic field quadratic over
\(\Q(\mathbf P(H))\),
\(\mathcal N\) is the normalization of \(\mathbf P(H)\) in that field.
For the icosahedron \(I_t\), let \(U_t\in X(E)\) be Hitchin's isotropic
three-plane.  The polarization identity and Hitchin's description give
\cite[\S\S2--5]{HitchinIcosahedron}
\[
 V_t=U_t^\perp,
\]
so \([p]\mapsto([p],U_t)\) defines a section
\(\mathbf P(V_t)\to\mathcal I_E\).  The two conjugate sections descend to
a morphism
\[
 Z^\nu\longrightarrow\mathcal I\longrightarrow\mathcal N
\]
over \(\mathbf P(H)\).  Over the finite etale neighborhood of \([xyz]\)
constructed below, the local comparison proves that this morphism carries
\(\nu^{-1}([xyz])=\Spec E\) isomorphically onto the golden incidence fibre.

\subsection{Proof of Theorem~\ref{thm:arithmetic-main}}

Hitchin realizes the Mukai--Umemura threefold as the zero locus of the
three skew forms on \(\operatorname{Gr}(3,H)\).  The top Chern number of
the dual universal bundle is \(2\), so a general \(f\) is orthogonal to
two isotropic three-planes \cite[\S\S4--5]{HitchinBundles}.  The incidence
variety is a projective bundle over the integral Mukai--Umemura threefold,
hence is integral; its generic function field is therefore a quadratic
extension.

\smallskip
\noindent\emph{The reduced branch cycle.}\par
Let \(\mathcal U\) be the universal rank-three bundle on
\(\operatorname{Gr}(3,H)\), put
\(\mathcal O(1)=\det\mathcal U^*\), and write \(h\) for the hyperplane
class on \(\mathbf P(H)\) and for its pullback to \(\mathcal I\).
The threefold \(X\) is the smooth zero locus of the rank-nine bundle
\(\bigwedge^2\mathcal U^*\otimes W\).  Since
\[
 K_{\operatorname{Gr}(3,7)}=\mathcal O(-7),\qquad
 \det(\bigwedge^2\mathcal U^*\otimes W)=\mathcal O(6),
\]
adjunction gives \(K_X=\mathcal O_X(-1)\).  Put
\[
 \mathcal F=\ker(H\otimes\mathcal O_X\longrightarrow\mathcal U^*).
\]
Then \(\det\mathcal F=\mathcal O_X(-1)\), and, for the bundle projection
\(p:\mathcal I\to X\), one has
\(\mathcal I=\mathbf P_X(\mathcal F)\), with its tautological
\(\mathcal O_{\mathcal I}(1)\) equal to
\(\pi^*\mathcal O_{\mathbf P(H)}(1)\).  The projective-bundle canonical
formula therefore gives
\[
 K_{\mathcal I}
 =\mathcal O_{\mathcal I}(-4)\otimes
   p^*(K_X\otimes\det\mathcal F^{-1})
 =\pi^*\mathcal O_{\mathbf P(H)}(-4).
\]
The determinant of \(d\pi\) is consequently a nonzero section of
\(K_{\mathcal I}\otimes\pi^*K_{\mathbf P(H)}^{-1}
 =\pi^*\mathcal O(3)\).  Its ramification divisor \(R\) has class \(3h\).
Because \(\pi\) has generic degree two, the projection formula gives the
branch cycle
\[
 \pi_*R=6h.                                                   \tag{2.1}
\]
Here a tamely ramified prime in a quadratic extension has ramification
index two and residue degree one, so each divisorial branch component occurs
once in \(\pi_*R\); components of \(R\) contracted to codimension at least
two do not contribute.

Hitchin proves that a general real point of \(J_0=0\) has exactly one
incidence configuration, whereas either side has respectively two or zero
real regular configurations
\cite[\S9, Proposition~10 and Theorem~11]{HitchinIcosahedron}.
At such a general boundary point \(\pi\) is quasi-finite, hence finite after
shrinking the target.  The source is smooth and the target regular of the
same dimension, so miracle flatness makes this local finite map flat of
degree two.  Its fibre has only one geometric point and therefore is
ramified.  The real smooth
points used here are Zariski dense in the irreducible sextic \(J_0=0\).
Thus this sextic is a component of the branch cycle.  Equation~(2.1) shows
that it exhausts that cycle, with multiplicity one.  In particular the
normal quadratic incidence field is branched exactly along the reduced
divisor \(J_0=0\).  This scheme-theoretic conclusion is the paper-local
ramification calculation; the cited real classification supplies the dense
set of boundary points, not the finite-cover assertion.

\smallskip
\noindent\emph{Square class from one fibre.}\par
The fibre over \([xyz]\) is the model case announced in the introduction: its
residue algebra will determine the constant twist of the global quadratic
field.  The first step below is an instance of a standard descent fact for quadratic
covers, which we isolate because it is the reason that a single arithmetic
fibre pins down a global twist.  Let \(X\) be a normal, geometrically integral
variety over a field \(k\) of characteristic not two, with
\(H^0(X,\mathcal O_X)=k\) and divisor class group free of two-torsion.
Let \(\mathcal L\) be a line bundle, let \(J\) be a nonzero section of
\(\mathcal L^{\otimes2}\) whose zero divisor is irreducible, and let \(M/k(X)\)
be a quadratic extension branched exactly along that divisor.  Then
\[
 M=k(X)\bigl(\sqrt{cJ}\bigr)
\]
for a square class \(c\in k^\times/k^{\times2}\) that is uniquely determined,
in the shorthand of the introduction.  Moreover, if \(x\in X(k)\) lies off the
branch divisor, then the normalization of \(X\) in \(M\) has residue algebra
\[
 k\bigl(\sqrt{c\,[J(x)]}\bigr)
\]
at \(x\), where \([J(x)]\in k^\times/k^{\times2}\) is the square class of the
value \(J(x)\) in the fibre \(\mathcal L_x^{\otimes2}\); this class is well
defined because changing a basis of \(\mathcal L_x\) rescales the value by a
square.  In particular, at a point where \(J\) takes a square value, a residue
algebra \(k(\sqrt d)\) gives \(c=d\).

Writing \(M=k(X)(\sqrt e)\), the odd
valuations of \(e\) occur exactly along the branch divisor, so
\(\operatorname{div}(e/j)\) is even for \(j=J/s^2\) with \(s\) a rational
section of \(\mathcal L\); the resulting divisor class is two-torsion, hence
trivial, so \(e=cjg^2\) with \(c\) a global unit, that is a constant.
Specialization at \(x\) then evaluates \(c\,[J(x)]\).  The point of the
statement is its last clause: one unramified rational point at which the
branch section takes a square value, and whose fibre is a quadratic field,
reconstructs every rational quadratic twist with the given branch divisor.
The comparison carried out below supplies exactly that square value, since
\(J_0(xyz)=(16/25)^2\).

\emph{Square class from the branch divisor.}
Write \(K=\Q(\mathbf P(H))\).  In characteristic zero every quadratic
extension of \(K\) has the form \(K(\sqrt d)\).  The prime divisors at
which \(d\) has odd valuation are exactly the branch divisors.  The
preceding ramification-cycle calculation proves that their reduced union is
the irreducible sextic \(J_0=0\).

Choose a nonzero rational section \(s\) of
\(\mathcal O_{\mathbf P(H)}(3)\) and set
\[
 j_s=\frac{J_0}{s^2}\in K^\times.
\]
Its odd valuations occur exactly along \(J_0=0\), because the divisor of
\(s^2\) is even.  Consequently
\[
 \operatorname{div}(d/j_s)=2D
\]
for a divisor \(D\) on \(\mathbf P(H)\).  The class of \(D\) is
two-torsion.  Since
\(\operatorname{Pic}(\mathbf P(H))=\mathbf Z\), \(D\) is principal.
Thus
\[
 d=cj_sg^2
\]
for some \(g\in K^\times\) and a unique
\(c\in\Q^\times/\Q^{\times2}\).  The incidence extension is therefore
\(K(\sqrt{cJ_0})\).

\emph{The local comparison at \(xyz\).}
Hitchin's classification gives exactly the two distinct configurations
\(I_t\) and \(I_{1-t}\) over \([xyz]\)
\cite[\S\S3, 7--9]{HitchinIcosahedron}.  Thus \(\pi\) is quasi-finite over
this point.  Because \(\pi\) is proper, there is a Zariski neighborhood
\(B\) of \([xyz]\) on which \(\pi\) is finite.  Let
\(\widetilde B\) be the normalization of \(B\) in the quadratic generic
field \(K(\sqrt{cJ_0})\).  The normal finite scheme
\(\pi^{-1}(B)\) is birational to \(\widetilde B\), hence is isomorphic to
it.  The branch-cycle calculation above identifies the branch locus with
the reduced divisor \(J_0=0\).  Since \(J_0(xyz)\ne0\), shrinking \(B\)
away from this divisor
makes the finite normalization etale.  The two distinct geometric points
therefore form the complete reduced scheme-theoretic fibre, whose residue
algebra is
\[
 E=\Q[t]/(t^2-t-1)=\Q(\sqrt5).
\]

Choose a local generator \(e\) of
\(\mathcal O_{\mathbf P(H)}(3)\) at \([xyz]\), and write
\(j_e=J_0/e^2\).  The normalization is locally
\(w^2=cj_e\).  Changing \(e\) multiplies \(j_e\) by a unit square.
Moreover, \(j_e(xyz)\) differs from the displayed value
\(J_0(xyz)=(16/25)^2\) by a rational square.  Specialization therefore
gives the residue algebra \(\Q(\sqrt c)\).  Comparison with the preceding
fibre proves
\(c=5\) in \(\Q^\times/\Q^{\times2}\).  This also explains why evaluating
the generic square class at the point is legitimate: the point lies in the
finite etale locus and the change from \(j_s\) to \(j_e\) is a unit square
there.

\emph{The Stein algebra.}
The quadratic field involution extends to the normalization \(\mathcal N\).
If \(f:\mathcal N\to\mathbf P(H)\), then, since \(2\) is invertible,
\[
 f_*\mathcal O_{\mathcal N}=\mathcal O\oplus\mathcal M
\]
is the decomposition into invariant and anti-invariant summands.  The
invariant summand is \(\mathcal O\): its relative spectrum is finite and
birational over the normal space \(\mathbf P(H)\).  Normality makes
\(f_*\mathcal O_{\mathcal N}\) an \(S_2\) module over the regular base, and
the direct summand \(\mathcal M\) is therefore a rank-one reflexive sheaf.
Every such sheaf on projective space is a line bundle, say
\(\mathcal M\simeq\mathcal O(-d)\).

Multiplication on the anti-invariant summand is a section of
\(\mathcal M^{-2}\simeq\mathcal O(2d)\).  Its zero divisor is the branch
divisor, which is the sextic \(J_0=0\); hence \(d=3\), and multiplication is
\(cJ_0\) for some \(c\in\Q^\times\).  The function-field calculation and
the golden fibre give \(c=5\) modulo a rational square.  Rescaling the chosen
isomorphism \(\mathcal M\simeq\mathcal O(-3)\) absorbs that square and gives
the exact algebra
\[
 f_*\mathcal O_{\mathcal N}
 \simeq\mathcal O\oplus\mathcal O(-3),
 \qquad z^2=5J_0.
\]
After Hitchin's scale for \(J_0\) and an isomorphism
\(\mathcal M\simeq\mathcal O(-3)\) giving multiplication \(5J_0\) are
fixed, the remaining choices differ by \(z\mapsto-z\), the deck involution.
Thus the equation used on the Clebsch chart below is an equation of finite
algebras, not only of their function fields.

\paragraph{Integral boundary.}
Choose a full lattice \(\Lambda\subset H\) and an integer \(a\ne0\) such
that \(J_{\mathbf Z}=a^2J_0\) has integral coefficients.  The algebra
\[
 \mathcal O_{\mathbf P(\Lambda)}
 \oplus\mathcal O_{\mathbf P(\Lambda)}(-3),
 \qquad z^2=5J_{\mathbf Z},
\]
is finite locally free of degree two over \(\mathbf Z\) and has the
function field of Theorem~\ref{thm:arithmetic-main}.  On
\(D(10J_{\mathbf Z})\) it is finite etale.  The primes \(2\) and \(5\)
are forced bad primes for an ordinary two-sheet interpretation: in
characteristic \(2\) the equation is inseparable, while in characteristic
\(5\) it is generically \(z^2=0\).

This abstract integral algebra does not by itself give an integral incidence
theorem.
The constructions in the cited Mukai--Umemura and Hitchin sources are over
\(\mathbf C\), and the available positive-characteristic comparison for
binary-sextic invariants is proved only in characteristic greater than five
\cite{BinaryForms}.  Spreading
out the rational incidence isomorphism gives some
\(\mathbf Z[1/N]\), with \(2,3,5\mid N\) after enlarging \(N\), but neither
the sources nor the present equations determine the other prime divisors
or prove that there are none.  Consequently every finite-field
icosahedral-incidence interpretation in this paper remains restricted to
an unspecified cofinite set of primes.  The formula for the abstract
quadratic algebra itself has no such extra exception.

The missing assertion is precise.  One must identify an integral harmonic
lattice and its three-skew-form Grassmannian zero locus with an integral
Mukai--Umemura model, prove flatness and normality, and show that Stein
formation and the chart sections commute with the required base changes.
Until those geometric steps are supplied, extra divisors of \(N\) are
evidence gaps rather than additional bad primes predicted by the quadratic
equation.
